\documentclass{article}
\usepackage{pathfinder-readable}

\title{Priced Certificates for Conditional-CVaR Obedience}

\begin{document}
\maketitle

\begin{abstract}
Q surveys guarantees and costs in learning-augmented systems, while P studies how finite-batch
conditional risk estimates can reverse an obedience decision. The thread asked whether a priced,
realised batch could instead certify that decision. For a fixed finite policy with access to losses
from deviations, the peers found a high-confidence reward certificate and a local bound on its
sampling and reward costs. The result needs a cycle margin for reliable acceptance and does not give
a system-wide cost guarantee, so the verifier left the thread at DRAFT.
\end{abstract}

\section{Introduction}

Q is a survey of algorithms that use predictions while retaining performance guarantees \cite{Q}. It
asks how to count the price of calling a predictor and when guarantees for separate parts of a
system can be combined. Such a combination needs conditions on each part and a common account of
system cost. Q itself gives no test for obedience to a prescribed action.

P considers a small incentive problem with fixed signals, actions, and rewards that depend on the
action but not on the signal \cite{P}. It values an action by the negative conditional value at risk
of its loss, or CVaR: the average loss in a specified upper tail. P shows that replacing true CVaR
by the \emph{expected} empirical CVaR of a finite batch can change whether a policy can be supported
by these rewards. It also gives a graph test when bounds on all action values are already known. P
leaves open a guarantee from one realised batch. The paper called Q within P is a different
mechanism-design paper, not the survey supplied here.

The peers pursued that gap. They asked whether samples could give a checkable reward for the true,
conditional-risk policy and what it would cost to obtain one. The question connects P's graph test
to Q's cost accounting. It concerns one fixed policy and its action-stage obedience. It does not
cover risk for a complete plan before the signal arrives or incentives to report a type.

\section{What was done}

The work began by separating P's two kinds of estimate. Its reversal concerns an expectation over
possible batches. A certificate must use the values calculated from the batch actually drawn. The
peers assumed a finite signal set \(S\), a set \(B\) of prescribed actions, and a policy \(a(s)\)
that prescribes an action at signal \(s\). Deviations to actions outside \(B\) were excluded. They
assumed access to independent draws within each action--signal cell, including cells for actions the
policy would not take. Each loss lies between zero and \(H\). These access and boundedness
assumptions are essential.

For each of the \(N=|S||B|\) cells, the peers took \(m\) draws. They called the upper-tail fraction
\(\beta\), with \(0<\beta\leq1\), and set the true utility \(U(b,s)\) to minus the conditional CVaR
of action \(b\) at signal \(s\). Its batch estimate is \(\widehat U(b,s)\). The CVaR variational
formula, the Dvoretzky--Kiefer--Wolfowitz inequality, and a union bound gave a simultaneous error
radius
\[
e_m=\frac{H}{\beta}
\sqrt{\frac{\log(2N/\delta)}{2m}}.
\]
Here \(\delta\) is the allowed failure probability. With probability at least \(1-\delta\), every
estimated utility differs from its true value by at most \(e_m\). Draws from different cells need
not be independent for this union bound. Existing work already studies empirical-CVaR concentration
and risk-arm identification \cite{cvarfixed,cvarother,cvarneurips}.

The peers then put these intervals into P's reward graph. An edge from prescribed action \(a\) to
another prescribed action \(b\) records the largest possible gain from switching at any signal that
prescribes \(a\):
\[
w_m(a,b)=
\max_{s:\,a(s)=a}
\bigl\{\widehat U(b,s)-\widehat U(a,s)+2e_m\bigr\}.
\]
They checked the graph for positive cycles and, where possible, solved for rewards. They next asked
whether the resulting rewards could be expensive even when the policy was feasible. Finally, they
tested what happens near a zero cycle margin, with only on-policy observations, and when this local
certificate is placed in Q's larger cost model. The detailed derivations and cross-checks are in
\url{../emmy/research.md} and \url{../ada/decision_certificate.md}.

\section{Results}

If the confidence graph has no positive directed cycle, there are rewards \(r(a)\) satisfying
\(r(a)-r(b)\geq w_m(a,b)\) for every edge. On the simultaneous coverage event, this same reward
obeys every true incentive inequality. The output can therefore include the estimates, the radius,
the graph check, and the reward vector as a verifiable certificate of exact true obedience, with
confidence at least \(1-\delta\). Conversely, a cycle whose empirical utility sum stays negative
even after adding \(2e_m\) for each edge certifies true infeasibility on that event. If neither
check succeeds, the answer is abstention. P supplies the deterministic graph argument; the peers
supplied the realised-sample interval that makes it usable here.

There is a sufficient sample size when cycles stay clear of zero. Suppose every nontrivial simple
true cycle has utility sum at least its number of edges times a positive margin \(\gamma\). Then the
feasible certificate succeeds on the coverage event if \(e_m<\gamma/4\). A true cycle whose sum is
at most minus its number of edges times \(\gamma\) similarly gives a negative certificate. It is
enough to choose
\[
m>\frac{8H^2}{\beta^2\gamma^2}
\log\frac{2N}{\delta}.
\]
If a loss draw costs \(\kappa_s\), the sample bill is \(\kappa_sNm\), plus any material cost of the
graph calculation. Q prices predictor invocations, so treating one draw as one invocation is a
further modelling choice.

The peers also bounded a local reward cost. Let \(n=|B|\). Let \(R^*\) be the smallest possible
maximum nonnegative reward that supports the policy using true utilities, and let \(R_m^*\) be the
corresponding minimum for the confidence graph. If every nontrivial simple true cycle has average
edge weight at most \(-\gamma\), and \(e_m<\gamma/4\), then on the coverage event
\[
R^*\leq R_m^*\leq R^*+4(n-1)e_m.
\]
The reason is that each confidence edge exceeds its true edge by at most \(4e_m\). The least
nonnegative reward at an action is the largest weight of a path starting there, including the empty
path. A largest path can be taken to have at most \(n-1\) edges. If a single reward payment counts
as resource cost, the resulting local upper bound is \(\kappa_sNm+R^*+4(n-1)e_m\). This compares
with an oracle that knows the true utilities; it is not a ratio against Q's offline system
benchmark. The peers' path argument and mutual check appear in \url{../emmy/research.md} and
\url{../ada/decision_certificate.md}.

\section{Objections and unresolved issues}

The first objection was statistical. In P's two-signal pattern, take a safe loss of \(1\) and a
risky loss of \(2X\), where \(X\) is a Bernoulli variable with success chance \(p\). At upper-tail
fraction \(1/2\), the true two-edge cycle is \(2-8p\) when \(p<1/2\). It changes sign across
\(p=1/4\), while every finite sample transcript remains possible on either side. The peers concluded
that a finite-sample rule cannot give a uniformly exact, zero-error answer at this boundary. A
nonabstaining rule with a fixed small error probability needs an expected number of informative
draws that grows as the inverse square of the distance from the boundary. Their change-of-measure
step uses a standard sequential testing argument \cite{klsource}. The margin promise and abstention
option resolve this objection for the stated certificate, not for unrestricted exact testing.

The second objection was access to deviations. The graph needs losses for an action even when the
policy never takes it at that signal. The peers gave two worlds with identical on-policy logs at
every sample size: prescribed losses are always \(3/2\), while unplayed deviations cost \(2\) in one
world and \(1\) in the other. One world has a positive obedience cycle and the other a negative one.
A conditional sampler, covered exploration, or a justified model of the missing losses is therefore
needed. The cost of obtaining such information must also be counted.

The third objection concerned money. A nonpositive cycle can still require an arbitrarily large
nonnegative reward. The peers addressed this locally by allowing an explicit reward cap and checking
it alongside the graph constraints. Their cap bound above applies only under its cycle margin and
cost convention. Earlier work also uses a positive value margin for incentive claims in a different
persuasion model \cite{persuasion}; this thread did not establish priority for its basic
concentration and margin ingredients.

Finally, confidence is not an expected system-cost bound. An expected-cost argument would need a
specified action on abstention and a cost bound on the coverage-failure event. It would also need
conditional coverage for every policy chosen after an adaptive history, plus a common downstream
objective and benchmark. The peers corrected an earlier cost claim that controlled only accepted
certificates: the cost bound must also cover abstention on the whole coverage event. None of those
system premises was established here.

\section{Why the thread stopped}

The verifier accepted the realised-batch certificate, the reward-cap comparison, and the boundary
obstruction as supported by the recorded arguments. It left the thread at DRAFT because the result
is a useful local connection, while the claimed extension to a system-wide competitive guarantee
remains unproved. The smallest restart is to specify a downstream objective and prediction-free
benchmark for the two-signal setting, say whether rewards count as costs, and give a feasible
abstention action with a cost bound. A conditional expected-cost comparison under those choices, or
a counterexample, would decide whether this certificate composes even in that minimal case.

\bibliographystyle{plainurl}
\bibliography{references}

\end{document}
